(1) 双体相对运动等效于约化质量 \(\mu\) 在中心势 \(-Gm_1m_2/r=-G\mu M/r\) 中运动。圆运动条件为 \(\mu r\omega^2=G\mu M/r^2\)，所以 \(\omega=\sqrt{GM/r^3}\)。

(2) 质心位置 \(m_1\boldsymbol r_1+m_2\boldsymbol r_2=0\)，相对位矢长度为 r，故两圆轨道半径 \(r_1=m_2r/M,\ r_2=m_1r/M\)，对应速率 \(v_i=\omega r_i\)，相对速度 \(v_{\rm rel}=\omega r=\sqrt{GM/r}\)。

(3) 总动能 \(K=\mu v_{\rm rel}^2/2=Gm_1m_2/(2r)\)，总势能 \(U=-Gm_1m_2/r\)，所以 \(E=-Gm_1m_2/(2r)\)。总绕质心角动量
\[
L=\mu r^2\omega=\mu\sqrt{GMr}.
\]
这是两个质点角动量之和，不可只用单个物体的质量。

(4) 逃逸边界的相对比能为0，需要 \(v_{\rm esc}=\sqrt{2GM/r}\)。初圆轨道相对速率 \(\sqrt{GM/r}\)，故纯切向的正速度增量
\[
\Delta v_{\rm min}=(\sqrt2-1)\sqrt{GM/r}.
\]
本结论是质心系相对速度增量，而非任意一颗物体的实验室单独脉冲。
